\documentclass[11pt]{article}
\usepackage[margin=0.9in]{geometry}
\usepackage{amsmath}
\usepackage{booktabs}
\usepackage{graphicx}
\usepackage{array}
\usepackage{longtable}
\usepackage{float}
\usepackage[T1]{fontenc}
\usepackage[colorlinks=true,linkcolor=blue,urlcolor=blue]{hyperref}
\graphicspath{{/home/runner/work/DataCenter/DataCenter/Presentation/scheduler_ramp_demo/}}
\setlength{\emergencystretch}{3em}
\setlength{\parskip}{0.35em}
\newcommand{\MW}{\,\mathrm{MW}}
\newcommand{\MWh}{\,\mathrm{MWh}}
\newcommand{\kWh}{\,\mathrm{kWh}}

\title{Small Changes in When Computing Starts\\
Can Make a Data Center Easier to Power\\[0.4em]
\large A plain-language guide to the scheduler ramp study}
\author{DataCenter repository}
\date{September 11, 2026}

\begin{document}
\maketitle

\begin{abstract}
This study asks a deliberately simple question: can a data center make its
power system's job easier by starting flexible computing work a few seconds
apart instead of starting everything at once? The experiment keeps the
requested work unchanged. It compares straightforward scheduling rules,
then gives each rule the same opportunity to use gas turbines and a
physically constrained battery. The power-management choice is the
least expensive feasible option from a small, explicitly listed set.
This is not a sophisticated scheduler, a global optimization proof, or a
validated design for a commercial data center.

The central distinction is between changing \emph{when} electricity is
needed and reducing \emph{how much} useful computing is done. The former can
reduce abrupt power changes without the latter. Any benefit must be
weighed against delayed work, storage losses, and the limits of the
power equipment. This report explains the physical ideas, the code,
the comparison method, the results, and the claims that the study
cannot support.
\end{abstract}

\tableofcontents
\clearpage

\section{The idea without any engineering background}

Imagine a building in which many elevators all start moving at once.
The motors suddenly ask the building's electricity supply for much more
power. Letting some elevators start a few seconds later can soften that
change, even if the elevators ultimately carry the same passengers.
The analogy is about timing, not about deliberately making useful work slower.

A data center has a similar opportunity when some of its work is flexible.
Its computers can start many tasks together, or spread those starts over
a short interval. The question is whether the delay is small enough to
be acceptable and the electrical benefit large enough to matter.

Three parts of the system have different jobs:
\begin{itemize}
\item \textbf{The scheduler} decides which waiting tasks may start.
It does not generate electricity.
\item \textbf{The gas turbines and generators} turn fuel into electrical
power. Their rotating machinery and fuel controls take time to respond.
\item \textbf{The battery} supplies or absorbs electricity quickly.
It cannot supply unlimited power or run indefinitely without recharging.
\end{itemize}

The useful coordination is simple: do not ask the slow equipment to
follow every sudden change if a small timing adjustment or a fast buffer
can help. But using a battery is not free, and delaying work is not free
from the customer's point of view.

\subsection{Power and energy are different}

\textbf{Power} is the rate of electricity use right now. \textbf{Energy}
is the amount used over time. A megawatt (MW) measures power; a
megawatt-hour (MWh) measures energy. One MWh is 1,000 kilowatt-hours (kWh).
For example, a 6 MW discharge lasting ten seconds delivers
\[
6\times\frac{10}{3600}=0.0167\MWh \simeq 16.7\kWh.
\]
This is a large power delivered for a short time, not a large energy store.

The battery needs both a sufficiently large power converter and enough
stored energy. A calculation that asks for several MW and only a few
tens of kWh is describing a short burst. It is not automatically a
practical battery pack specification.

\subsection{What is a ramp?}

Here, a ramp means the change in power between consecutive one-second
samples. A jump from 10 MW to 16 MW is a 6 MW/s one-second change.
Spreading the same increase over twelve seconds can make the changes
much smaller.

The qualifier ``one-second'' matters. The turbine simulation receives
loads that are held constant between updates and jump at the updates.
The study does not establish an upper bound on the instantaneous or
millisecond-scale change in a real computer's power.

\subsection{Why frequency appears in the results}

The generators rotate. If electrical demand rises before mechanical
power has caught up, energy comes temporarily from that rotation and
the machines slow down. Electrical frequency measures this rotation
indirectly. The model's normal reference is 60 Hz.

For this study, a coordinated operating choice must stay between
59.4 and 60.6 Hz at the sampled times. A sampled crossing below
57.8 Hz ends the published generator-only screening trajectory.
These are study thresholds, not universal equipment guarantees.
The model does not simulate what every device does after a plant trips.

\section{The question, and the meaning of ``optimal''}

The intended claim is:
\begin{quote}
Small, explainable adjustments to the timing of flexible work can improve
the tradeoff between computing service and power-system operation.
\end{quote}

That is narrower than claiming that all assets operate at a global optimum.
In this experiment:
\begin{enumerate}
\item The two generators are already online and share the load equally.
Their number, on/off status, and individual economic dispatch are not optimized.
\item The same installed battery is used when comparing schedules.
Separate hardware sensitivity cases change that assumption explicitly.
\item Every schedule is offered the same five power-management choices.
\item Choices that violate the engineering checks are rejected.
\item The remaining choice with the lowest modeled fuel and storage-use
cost is selected.
\end{enumerate}

Thus ``best'' means \textbf{best feasible tested choice}, not best possible
control law. Searching a small list of settings keeps the result transparent:
the reader can see both the alternatives and why an alternative was rejected.

\subsection{The cost objective}

The illustrative operating cost is
\[
C = 0.25 F + 20 Q,
\]
where $F$ is total fuel consumed by the fleet in kg and $Q$ is the sum
of battery charging and discharging energy in MWh. The assumed prices
are \$0.25/kg of fuel and \$20/MWh of battery throughput.

The battery term is a simple use-cost assumption, not a chemistry-specific
aging calculation. The objective omits equipment purchase, financing,
maintenance, emissions prices, and the monetary value of job delay.
Delay is instead treated as a constraint. Small dollar differences in this
short synthetic experiment must not be extrapolated directly to annual savings.

\section{What the computer workload represents}

A \textbf{node} is one modeled server. Each Frontier node represents four
GPUs, specialized processors often used for AI work. A \textbf{cohort}
is a group of nodes that starts together and runs for the requested duration.

There are nine batch jobs, arriving in three waves. The first wave requests
2,800 nodes, the second 2,700, and the third 2,400. In the synchronized
case the first two waves run for 240 seconds per cohort and the last
for 180 seconds. Arrivals occur at 120, 360, and 600 seconds.
The simulation begins with a quiet background before the first arrival.

The total is 7,900 node allocations across time, not 7,900 additional
physical servers. The synchronized workload contains
\[
2800(240)+2700(240)+2400(180)=1{,}752{,}000
\]
useful node-seconds. A node-second means one node doing the modeled useful
work for one second. A cohort that starts later also finishes later;
its requested execution time is not shortened.

\subsection{What flexibility is assumed}

The batch work is assumed to consist of independently startable tasks.
This is not a model of tightly synchronized distributed training,
where many workers may have to begin together. It is also not a claim
that individual online inference requests can wait several seconds.
The constant background represents service that should not be delayed.
The code also supports protected jobs that must start in full immediately;
it rejects a scenario if capacity cannot meet that requirement.

The work proxy proves accounting consistency within this model.
It does not measure model-training convergence, tokens generated,
customer response times, or AI quality. Communication overhead,
startup energy, and changes in task speed with power are outside the model.

\subsection{Three understandable scenarios}

\begin{description}
\item[Synchronized.] Equal durations within each wave; 1,600 background
nodes. Replacement work is conveniently available near earlier completions.
This is an intentionally favorable introductory example.
\item[Irregular.] The same job sizes and arrivals, but durations of
190, 270, 310, 175, 265, 205, 180, 230, and 150 seconds.
Completions are less aligned; total useful work is 1,716,500 node-seconds.
\item[Busy.] The synchronized jobs with 6,800 background nodes. Only
2,672 of the 9,472 available nodes remain for batch jobs, so even
immediate admission must queue some work.
\end{description}

The busy example is important: not every observed wait is caused by the
smoothing rule. Some waiting would occur anyway because the facility is full.

\section{The scheduling rules}

All rules use the same arrival order, with job ID as a deterministic tie-breaker.
They are not learning algorithms, priority optimizers, or deadline predictors.

\begin{enumerate}
\item \textbf{Immediate:} start all waiting work that fits.
\item \textbf{512-node wave:} start no more than 512 flexible nodes each second.
\item \textbf{256-node wave:} use the same rule with 256.
\item \textbf{226-node wave:} use the same rule with 226.
\item \textbf{Reuse plus 226:} start useful replacement work on nodes
whose previous work just finished, then permit up to 226 additional
active nodes. Count capacity before admitting any work.
\end{enumerate}

Why 226? The calibration adds about 0.002209 MW for each node that goes
from idle to full activity. The whole-node budget below 0.5 MW/s is
\[
\left\lfloor\frac{0.5}{0.002209}\right\rfloor=226.
\]
The matched 226-node wave separates the effect of a smaller admission
budget from the effect of the reuse rule.

All strategies can physically reuse nodes. The difference is that fixed
waves count replacement starts against their start budget, whereas the
reuse rule can admit them in addition to newly powered activity.

\subsection{No artificial full-power holding}

When work finishes, it stops using the active-power increment.
If there is no useful replacement work, the node returns to idle.
The model does not burn extra electricity merely to make a plot smoother.
Idle equipment still consumes the calibrated idle power.

This makes shutdowns honest. New starts can be paced, but previously
started tasks can finish together. Reuse plus 226 is therefore not a
guarantee of a symmetric $\pm0.5$ MW/s facility ramp.
It can also create overlap between old and new waves and perform worse
than the simpler fixed-wave rule.

\section{From active nodes to electrical demand}

The Frontier component configuration gives 9,472 available nodes,
about 7.307 MW of facility idle demand, and about 2.209 kW of additional
power per fully active node. The calculation includes the modeled CPU,
GPU, memory, networking, conversion losses, and cooling-distribution
components.

At each second the facility demand is
\[
P_{\mathrm{facility}}=P_{\mathrm{idle}}+
N_{\mathrm{active}}\Delta P_{\mathrm{node}}.
\]
The active count includes the fixed background and useful batch work.
In the synchronized case, 1,600 background nodes bring the pre-arrival
demand to about 10.841 MW. Adding the first 2,800 batch nodes brings
the immediate peak to about 17.025 MW.

This is a fixed component-based screening model. It does not model
changing weather, dynamic cooling, networking bottlenecks, or actual
GPU utilization within a task. It should not be presented as a
measurement of a present-day commercial AI campus.

\section{Three distinct power calculations}

\subsection{Generator-only reference}

The facility demand alone is supplied to the GGOV1 multishaft model.
GGOV1 represents fuel-control behavior; the multishaft extension represents
the gas-generator and power-turbine rotational response.
The two units are represented as one equal-sharing dynamic equivalent.
Each has a 22 MW turbine power base and 23 MVA generator base.

For this study, electrical demand is not reduced when frequency falls.
This prevents an apparently successful result from relying on computing
work silently consuming less power. It also changes the reference response
relative to the older study, which used frequency-sensitive load damping.
Old frequency numbers cannot be combined with the revised results.

The unbuffered run is a diagnostic, not necessarily an acceptable operating
choice. If it crosses the protection threshold, its published trace ends
at that first sampled crossing and its full-study fuel total is left missing.

\subsection{Ideal fast-buffer estimate}

For continuity with the original demonstration, the code also calculates
a lossless, unlimited ten-second smoothing split. The slowly changing
target is assigned to the generator and the fast difference to an ideal buffer.
This answers a useful preliminary question: how much short-duration
power exchange would this particular filter request?

It does not enforce a real battery's energy, efficiency, or inverter limits.
Its energy swing is not the installed energy rating. Ratios of hundreds
of times the energy rating per hour indicate a very short pulse requirement,
not an ordinary battery configuration.

\subsection{Constrained generator--battery operation}

The coordinated calculation instead starts with actual installed ratings:
\begin{center}
\begin{tabular}{ll}
\toprule
Assumption & Default value\\
\midrule
Battery maximum power, either direction & 6 MW\\
Battery nominal stored energy & 1.5 MWh\\
Allowed stored charge & 10--90\%\\
Initial stored charge & 50\%\\
Charging efficiency & 95\%\\
Discharging efficiency & 95\%\\
Generator reserved headroom & 2 MW\\
Maximum one-second generator load change & 0.5 MW/s\\
Final quiet energy-recovery period & 120 s\\
\bottomrule
\end{tabular}
\end{center}

The generator headroom is measured below the model's continuous ceiling.
It is not a demonstration that the plant can survive losing one generator.
The assumed inverter follows its one-second command perfectly within its
power limit; switching dynamics and voltage control are not represented.

Positive battery power supplies the facility. Negative power charges the
battery. For every dispatch interval,
\[
P_{\mathrm{generator,electrical}}+P_{\mathrm{battery}}
=P_{\mathrm{facility}}.
\]
This is an electrical allocation, not an assumption that the turbine's
mechanical power instantly equals its electrical load. The turbine model
receives the \emph{remaining electrical demand}. Temporary mechanical
imbalance changes the rotating energy and frequency.

Battery energy changes according to
\[
E_{k+1}=E_k+
\frac{\eta_c P_{\mathrm{charge},k}
-P_{\mathrm{discharge},k}/\eta_d}{3600}.
\]
Powers are in MW, energies in MWh, and each interval lasts one second.
Charging and discharging losses mean that more electricity must be put back
than was delivered to the facility.

Storage commands are clipped to physical limits. Any resulting generator
ramp or headroom violation is reported rather than hiding the unmet demand.
A saturated battery is not automatically an infeasible case: it is infeasible
if the resulting system fails the actual engineering checks.

\subsection{The five tested operating choices}

The choices are no storage action and smoothing times of 2, 5, 10, and
20 seconds. In the latter cases a target follows demand gradually.
The battery attempts to supply the difference.

Near the end of the study, a scheduled quiet period restores the initial
stored energy. The recovery period knows the experiment's endpoint.
This is not a forecast-free controller for an endlessly running facility.
Candidates that cannot restore the charge are rejected.

The energy accounting uses the intervals $[0,T)$.
The row at $T$ is an endpoint, not an extra second of fuel demand or battery
use. Every schedule within a scenario is padded to the same endpoint,
at least 240 seconds after the latest completion. This avoids counting
extra background operation for only the slower schedules.

\section{Reading the measurements correctly}

\begin{description}
\item[Useful work.] The integral of useful batch nodes over time.
It must match across schedules within each scenario.
\item[Full-start delay.] The time from a job's arrival until its last cohort
starts. This is not the delay before its first cohort starts.
\item[Additional delay.] A job's full-start time minus that same job's
full-start time under immediate admission in the same scenario.
\item[Completion extension.] How much later the final job finishes than
under immediate admission. It need not equal the largest per-job delay.
\item[Makespan.] Time from the first arrival to the final completion.
The separate last-completion field is a timestamp from simulation start.
\item[Maximum ramp.] The largest absolute one-second power change.
\item[Event percentile.] A percentile calculated only over nonzero changes.
A whole-trace 99th percentile can misleadingly be zero when changes are rare.
\item[Feasible policy.] A tested power setting that passes every required
limit. Missing best cost means none passed; it does not mean free operation.
\end{description}

The sensitivity analysis accepts the best tested schedule and power setting
whose \emph{maximum additional per-job delay} fits an allowance of 0, 5, or
15 seconds. It does not impose a 15-second bound on all capacity waiting.
For example, the busy immediate baseline can already have a long queue.

\section{Results of the completed study}
\label{sec:results}

The numerical comparisons in this section use regenerated output, not the
older four-strategy results. Generator-only values and coordinated values
must not be interchanged. Each hardware comparison uses a fixed installed
battery; comparisons between hardware sizes are sensitivity studies, not
capital-cost optimization.

\subsection{The introductory case: a small delay, a much gentler power change}

Table~\ref{tab:scheduling} describes the scheduling side before a battery
does anything. Each row completes exactly 1,752,000 useful node-seconds
and consumes 4.1765 MWh of modeled facility energy over the common
1,030-second accounting window.

\begin{table}[H]
\centering
\small
\begin{tabular}{lrrr}
\toprule
Rule & Largest change & Largest added & Final completion\\
 & (MW/s) & job delay (s) & extension (s)\\
\midrule
Immediate & 6.185 & 0 & 0\\
512-node wave & 1.131 & 5 & 4\\
256-node wave & 0.565 & 10 & 9\\
226-node wave & 0.499 & 12 & 10\\
Reuse plus 226 & 1.498 & 12 & 3\\
\bottomrule
\end{tabular}
\caption{Synchronized scheduling results. The delay is the largest extra
full-start wait of any job, not the delay of every task.}
\label{tab:scheduling}
\end{table}

The 226-node wave reduces the largest one-second change by about
\textbf{92\%}, while the final job finishes ten seconds later. The
immediate case finishes at simulation second 780; the 226-node wave
finishes at second 790. Measured from the first arrival, these durations
are 660 and 670 seconds, respectively.

This is not an energy-saving trick: facility energy is the same.
The electrical benefit comes from changing \emph{when} the work runs.
The 256-node rule already achieves roughly a 91\% ramp reduction.
That supports beginning with a simple rule rather than claiming that
special scheduling intelligence is required.

\subsection{What the same generators and battery do with those schedules}

Table~\ref{tab:operation} gives the least costly feasible tested power
setting for each schedule using the same 6 MW / 1.5 MWh installed battery.
The smoothing time is a controller setting, not the amount of job delay.

\begin{table}[H]
\centering
\small
\begin{tabular}{lrrrr}
\toprule
Rule & Smoothing & Storage peak & Fuel & Cost\\
 & time (s) & (MW) & (kg) & (\$)\\
\midrule
Immediate & 20 & 5.890 & 1264.43 & 317.56\\
512-node wave & 10 & 4.379 & 1264.27 & 316.79\\
256-node wave & 5 & 2.417 & 1264.18 & 316.41\\
226-node wave & 0 & 0.000 & 1264.10 & 316.02\\
Reuse plus 226 & 10 & 4.121 & 1264.33 & 316.98\\
\bottomrule
\end{tabular}
\caption{Synchronized coordinated operation, including battery losses
and restoration of its starting charge. Zero smoothing means no battery use.}
\label{tab:operation}
\end{table}

Without battery support, immediate admission crosses the selected
57.8 Hz protection threshold. Its coordinated 20-second policy stays
within the required frequency band and is therefore a legitimate
baseline for cost comparison. Comparing a successful schedule against
the fuel consumed by an interrupted baseline would not be fair.

In this particular synchronized example, the 226-node wave passes the
checks with no battery action. Its frequency spans about
59.472--60.178 Hz. Immediate admission with its selected battery policy
spans about 59.663--60.000 Hz. Thus the simpler schedule uses less
storage, but does \emph{not} deliver the best frequency margin.
Both choices pass the stated checks.

The modeled operating-cost reduction is only about \$1.53, or
0.48\%, over this short window. Fuel falls by about 0.34 kg, or
0.027\%. Most of the cost difference comes from the assumed battery-use
charge, not a large improvement in turbine efficiency. Immediate
admission exchanges about 72.43 kWh of total charging-plus-discharging
energy and loses about 3.71 kWh in storage; the selected 226-node
case exchanges none.

The separate ideal ten-second calculation asks for 5.622 MW under
immediate admission and 3.401 MW under the 226-node wave, a reduction
of about 40\%. Both ideal energy swings remain 17.18 kWh.
These are \emph{not} the actual storage peaks in
Table~\ref{tab:operation}, which uses different, constrained
power-management choices.

\subsection{Why the reuse rule is not the headline winner}

The matched comparison reveals a real tradeoff. Reuse plus 226 finishes
the last job only three seconds later than immediate admission,
instead of ten seconds later for the fixed 226-node wave.
But the reuse rule also admits additional work while earlier cohorts
are still running, which can concentrate later completions.
Its largest one-second change is about three times that of the
fixed rule, and it still needs storage action in its chosen policy.

This is a useful negative result, not a failure to hide.
Reusing capacity quickly may favor completion time, while pacing all
starts may favor electrical smoothness. Neither rule should be called
universally superior.

\subsection{Irregular completions: smaller ramps are not enough by themselves}

For unequal task durations, all five schedules preserve 1,716,500 useful
node-seconds and 4.2963 MWh of facility energy over 1,077 seconds.
Immediate admission selects 20-second smoothing at about \$332.42.
The 512-node wave selects ten seconds at \$331.03. The 256-node
and 226-node waves select five seconds at about \$330.10.

Unlike the synchronized example, the 226-node wave is \emph{not}
acceptable without battery support. Its generator-only frequency falls
to about 59.282 Hz, even though its largest one-second facility change
is only 0.499 MW/s. Multiple changes close together can cause a worse
dynamic response than a single isolated change with the same ramp.

The selected 226-node coordinated case peaks at 2.216 MW of battery
power and reaches about 59.407 Hz at its lowest point. That is only
about 0.007 Hz above the study's lower limit: it is a narrow screening
margin, not a robust deployment guarantee. The final completion is
seven seconds later than immediate admission, and the largest
additional full-start wait is twelve seconds.

Here the reuse and fixed-226 rules produce identical schedules.
The 256-node and 226-node costs differ by less than one tenth of a
cent. Treat them as effectively tied for practical interpretation,
not as evidence that a specific integer budget is economically optimal.

\subsection{Busy cluster: distinguish capacity waiting from intentional delay}

The busy case reaches all 9,472 available nodes. Every strategy has a
maximum full-start wait of 240 seconds because some requested work
cannot fit until earlier work completes. It would be misleading to
describe that four-minute wait as a cost imposed by ramp control.

Relative to the same jobs under immediate admission, the additional
maximum waits are four seconds for the 512-node rule, nine for the
256-node rule, and eleven for either 226-node rule. All five strategies
finish their last job at the same time, second 840. The capacity
bottleneck, rather than the short staggering interval, controls that
last completion.

All preserve 1,752,000 useful node-seconds and 7.7728 MWh of facility
energy in the common 1,080-second window. The best tested default-hardware
cost falls from \$433.29 for immediate admission to \$431.79 for the
fixed-226 rule. Their raw maximum ramps are 5.902 and 0.499 MW/s.
Immediate admission uses a 5.621 MW storage peak; the fixed-226 and
identical reuse case use no storage action. The different total cost
from the synchronized case reflects a different background load and
accounting window, not a direct measure of one scenario's superiority.

\subsection{How much scheduling flexibility changes the feasible choices}

Table~\ref{tab:allowances} summarizes the default-hardware comparison.
An allowance limits the largest extra per-job full-start delay relative
to immediate admission, not the duration of an entire job.

\begin{table}[H]
\centering
\small
\begin{tabular}{lrrr}
\toprule
Scenario & No extra delay & At most 5 s & At most 15 s\\
\midrule
Synchronized & Immediate / \$317.56 & 512 / \$316.79 & 226 / \$316.02\\
Irregular & Immediate / \$332.42 & 512 / \$331.03 & 226 / \$330.10\\
Busy & Immediate / \$433.29 & 512 / \$432.54 & 226 / \$431.79\\
\bottomrule
\end{tabular}
\caption{Best feasible tested schedule and operating cost with the
6 MW / 1.5 MWh hardware. ``226'' includes the identical reuse alternative
in the irregular and busy cases.}
\label{tab:allowances}
\end{table}

With either smaller hardware option (2 MW / 0.5 MWh or 4 MW / 1 MWh),
no tested combination passes all checks under the zero- or five-second
allowance in any scenario. Under the fifteen-second allowance,
at least one tested combination is feasible in all three scenarios.
This is the clearest system-level benefit: limited scheduling flexibility
changes what a given set of power assets can support.

These statements concern the tested policies only. A different controller
might make an excluded combination feasible. Nor does this establish
an optimal battery purchase: purchase costs, contingency operation, and
other storage duties are outside the experiment.

In the irregular case, the 2 MW battery saturates for eleven seconds
under the selected fixed-226 policy, yet the resulting generator load
and frequency still pass the checks. Its cost is about \$330.08,
slightly below the larger-battery result because it performs slightly
less storage exchange. This illustrates why clipping is modeled
explicitly and why bigger hardware is not automatically the cheapest
operating choice under a fixed family of control rules.

\subsection{How to read the figures}

\begin{figure}[H]
\centering
\includegraphics[width=0.96\linewidth]{01_scheduler_power_ramp.png}
\caption{Synchronized workload: active GPUs, facility demand, and one-second
power changes. Look first at the opening rise and the final drop. Spreading
starts changes the shape, not the amount of requested useful work.}
\end{figure}

\begin{figure}[H]
\centering
\includegraphics[width=0.96\linewidth]{02_turbine_response.png}
\caption{Generator-only screening. This is not the coordinated battery case.
A trace that ends at a protection crossing is an unacceptable reference,
not evidence of continued operation. Exhaust temperature is an estimate.}
\end{figure}

\begin{figure}[H]
\centering
\includegraphics[width=0.96\linewidth]{03_strategy_tradeoffs.png}
\caption{Different questions in one view: raw facility ramps, ideal-buffer
power, unbuffered frequency versus delay, and best feasible coordinated
operating cost. The ideal-buffer panel is not an installed-storage rating.
The cost panel uses the default fixed hardware.}
\end{figure}

\begin{figure}[H]
\centering
\includegraphics[width=0.96\linewidth]{04_fast_buffer_comparison.png}
\caption{Coordinated immediate and fixed-226 operation with the same
installed hardware. Electrical generator load is distinct from mechanical
turbine power. Battery charge is restored at the end, so stored energy
does not subsidize the apparent operating cost.}
\end{figure}

\section{How the code produces the study}

The main entry point is
\path{/home/runner/work/DataCenter/DataCenter/tools/workload/scheduler_ramp_demo.py}.
Storage accounting is in
\path{/home/runner/work/DataCenter/DataCenter/tools/workload/power_asset_dispatch.py}.
The code follows the same order as the explanation:
\begin{enumerate}
\item \texttt{load\_frontier\_power\_calibration} reads component and topology
values and derives idle demand and the active-node increment.
\item \texttt{demo\_workload} constructs one of the three transparent scenarios.
\item \texttt{simulate\_schedule} advances one second at a time. It retires
finished cohorts, registers new arrivals, admits protected and flexible work,
and records active nodes, queues, and job times.
\item \texttt{align\_schedule\_windows} pads completed schedules with background
activity so energy and cost share a common observation period.
\item \texttt{add\_buffer\_columns} retains the separate, unlimited ideal estimate.
\item \texttt{dispatch\_storage} applies power and stored-energy constraints,
tracks losses, allocates residual generator demand, and restores charge.
\item \texttt{run\_turbine} replays the generator's electrical demand through
the dynamic model. Fuel is reported for the whole fleet.
\item \texttt{compare\_asset\_policies} rejects unacceptable candidates and
chooses the least costly feasible tested setting.
\item \texttt{main} repeats the experiment for hardware sizes and scenarios,
applies the delay allowances, and writes tables, traces, figures, and a manifest.
\end{enumerate}

No candidate scheduler predicts a turbine state or solves a complicated
job-placement optimization problem. The separation is intentional:
simple computing decisions are evaluated fairly against the same
power-management capabilities.

\subsection{Output guide}

The output root is
\path{/home/runner/work/DataCenter/DataCenter/Presentation/scheduler_ramp_demo/}.
The synchronized case uses the root; other cases use the
\texttt{irregular} and \texttt{busy} subdirectories.

\begin{description}
\item[\texttt{strategy\_summary.csv}] One row per strategy.
Unprefixed frequency and temperature columns are generator-only results;
\texttt{best\_} columns describe coordinated operation.
\item[\texttt{dispatch\_candidates.csv}] All tested hardware/settings, costs,
losses, feasibility decisions, and rejection reasons. Missing dynamic
results can mean a candidate failed a simpler check before dynamic replay.
\item[\texttt{job\_delays.csv}] Individual job arrival, start, and finish times.
\item[Parquet traces.] Scheduler and turbine trajectories, plus the chosen
coordinated dispatch and its turbine replay where a feasible choice exists.
For a coordinated dispatch, row number is time in seconds.
\item[\texttt{scenario\_comparison.csv}] Combined scenario summaries.
\item[\texttt{sensitivity.csv}] Best feasible choices for each hardware
and additional-delay allowance.
\item[\texttt{run\_manifest.json}] Arguments, hardware assumptions, package
versions, and source/configuration hashes for the run.
\end{description}

\section{Verification, reproducibility, and numerical limits}

The existing tests check that work is preserved, no capacity is exceeded,
protected jobs start immediately or fail explicitly, and completed nodes
are not artificially held at full power. They also check physical storage
power and energy balance, losses, charge restoration, saturation,
infeasible cases, fleet scaling, protection handling, and finite-cost
selection.

Two additional regression checks protect reproducibility and accounting.
A cached turbine response is only reused for the same fleet, time grid,
sampling interval, and electrical load. The output includes the exact
study endpoint even when the chosen sample interval does not divide
the duration evenly; the final shorter interval still contributes to fuel
accounting. Temperature-change rates use actual time differences.
The existing test suite passes all 65 tests, and the separate smoke test passes.

The documented Python entry point runs all three scenarios and storage
sensitivities by default. The existing Pixi task
\texttt{pixi run scheduler-ramp-demo} invokes the same entry point when
used in the configured workspace. The workspace currently targets
Apple Silicon; Python dependencies were installed separately for the Linux
run described here, and the manifest records the actual package versions.
No project dependency was changed for that run.

The report can be built with the repository's Tectonic dependency.
For this validation run, Tectonic's default bundle host could not be
resolved, so the report was compiled with pdfLaTeX from TeX Live instead.
Two passes resolve the table of contents and cross-references.

The dynamic output interval is 0.05 seconds by default.
Frequency acceptance is based on sampled results, not a proof that no
between-sample excursion exists. The simulation is a deterministic
comparison, not a statistical estimate from many real workload days.
Close rankings should be treated cautiously rather than advertised with
more precision than the model justifies.

As a focused numerical check, the irregular case with default hardware
was repeated at a 0.025-second output interval. All five selected
power-policy settings remained unchanged. Selected frequency extrema
changed by less than $0.000006$ Hz and operating costs by less than
\$0.000001. This checks output sampling, not uncertainty in physical
parameters or numerical convergence of every model equation.

\section{Temperature, shaft stress, and what not to claim}

The exhaust estimate maps an existing filtered fuel-temperature signal
through a ThermoPower part-load curve and shifts it to a documented
LM2500 full-load exhaust point of 791 K. This is not blade metal temperature.
The part-load shape and transient response are not calibrated to this
particular generator installation.

A lower peak rate of exhaust-temperature change may suggest a gentler
transient. It does not establish reduced metal stress, fewer damaging
cycles, or a known life extension. A longer sequence of smaller changes
can have different total variation from one sharp change.

The code retains optional torque and fatigue helpers, but those are not
part of the default headline study. The underlying gas-path coupling,
shaft dimensions, materials, and thermal states need appropriate validation
before numerical fatigue damage can support equipment-life claims.

\section{Conclusions and appropriate next steps}

The important result is not that the scheduler is clever. The useful
question is whether a modest amount of timing flexibility changes what
the \emph{same power equipment} can do while respecting the computing
service requirement.

The regenerated examples support a concrete but limited conclusion:
staggering independent tasks by up to twelve seconds substantially
reduces abrupt facility power changes, and can make smaller tested
storage configurations feasible. This does not require a clever
scheduler. The matched reuse rule provides no consistent additional
electrical benefit, and in the synchronized example is worse than
simple fixed waves.

Facility energy and requested useful work are unchanged. Modeled
operating-cost improvements are small, and largely reflect reduced
storage use rather than large fuel savings. Irregular completions still
require battery support and leave a narrow frequency margin. The
evidence therefore supports coordinating simple scheduling flexibility
with power assets, not promising universal energy savings or global
optimality.

A defensible interpretation distinguishes three outcomes:
\begin{enumerate}
\item \textbf{Electrical smoothing:} the facility presents smaller abrupt
changes to its power supply.
\item \textbf{Operational feasibility:} a combination of scheduling and
storage passes the selected engineering limits.
\item \textbf{Economic preference:} among feasible tested alternatives,
one has a lower modeled fuel-plus-storage-use cost.
\end{enumerate}
None of these automatically proves the other two. Nor does any establish
a global optimum, improved AI throughput, or a smaller commercially
viable battery without further evidence.

For a real deployment, the next steps are to validate which workloads
can be delayed, calibrate the plant and storage models, use realistic
long-duration traces, and include reserve requirements, inverter and
cooling dynamics, equipment costs, and battery lifetime.
The simple scheduling rules should remain the starting point; additional
complexity should be introduced only if a measured benefit justifies it.

\appendix
\section{Glossary}
\begin{description}
\item[GPU.] A processor suited to highly parallel calculations.
\item[Node.] One modeled server containing CPUs, GPUs, and supporting hardware.
\item[Cohort.] Nodes admitted together for the same requested execution time.
\item[Scheduler.] The rule deciding when waiting computing work may start.
\item[Islanded.] Operating without a large external grid supplying the balance.
\item[Generator.] The electrical machine driven by the turbine.
\item[Governor.] The control that adjusts fuel in response to operating conditions.
\item[BESS.] Battery energy storage system, including its power converter.
\item[State of charge.] Stored energy as a fraction of nominal battery energy.
\item[Headroom.] Capacity deliberately left unused to provide an operating margin.
\item[Throughput.] Here, total charging plus discharging energy; it does not mean
computing throughput.
\item[Screening model.] A simplified model for exploring possibilities before
detailed design and validation.
\item[Feasible.] Passing the stated checks, not necessarily meeting every
requirement of a real installation.
\end{description}

\end{document}
